%------------------------------------------------------------------------------%
\documentclass[fleqn,utf8,aspectratio=169,12pt]{beamer}

\usepackage{calculo_varias_variaveis-1}

\title{Revisão de Limites}

%------------------------------------------------------------------------------%
\begin{document}

\addtitle

%------------------------------------------------------------------------------%
\section{Limites}

%------------------------------------------------------------------------------%
\begin{frame}{Limite de uma Função}
  Seja $f(x)$ definida em um intervalo aberto em torno de $a$,

  \pause
  exceto, possivelmente, no próprio $a$

  \pause
  O limite de $f(x)$ quando $x$ tende para $a$ é o número $L$
  \[
    \lim_{x\to a} f(x) = L
  \]
  \pause
  se, para cada \(\epsilon > O\), existir \(\delta > O\), tais que
  \pause
  \[
    0 < \abs{x-a} < \delta \qquad\Rightarrow\qquad \abs{f(x)-L} < \epsilon
  \]
\end{frame}



%------------------------------------------------------------------------------%
\begin{frame}{Limites de algumas funções}
  \( \lim_{x\to a} c = c \)

  \pause
  \vspace{0.5\baselineskip}
  \( \lim_{x\to a} x = a \)

  \pause
  \vspace{0.5\baselineskip}
  \( \lim_{x\to a} x^n = a^n \)
  \tabto{35mm}
  $n$ número inteiro positivo

  \pause
  \vspace{0.5\baselineskip}
  \( \lim_{x\to a} \sqrt[n]{x} = \sqrt[n]{a} \)
  \tabto{35mm}
  $n$ número inteiro positivo, se $n$ for par precisamos impor $a\geq 0$

  \pause
  \vspace{0.5\baselineskip}
  \( \lim_{x\to a} p(x) = p(a) \)
  \tabto{35mm}
  onde $p(x)$ é um polinômio em $x$
\end{frame}

%------------------------------------------------------------------------------%
\begin{frame}{Propriedades de Limites}
  Se os limites \(\lim_{x\to a}f(x)\) e  \(\lim_{x\to a}g(x)\) existirem, então:\\[4mm]
  \pause
  \begin{enumerate}[<+->]
    \item
      \(
        \lim_{x\to a} \left(f(x)+g(x)\right)
        = \lim_{x\to a}f(x) + \lim_{x\to a}g(x)
      \)
      \hfill Soma
      \\[5mm]

    \item
      \(
        \lim_{x\to a} \left(f(x)-g(x)\right)
        = \lim_{x\to a}f(x) - \lim_{x\to a}g(x)
      \)
      \hfill Diferença
      \\[5mm]

    \item
      \(
        \lim_{x\to a} \left(c f(x)\right)
        = c\lim_{x\to a} f(x),\quad \forall\; c\in\R
      \)
      \hfill Multiplicação por escalar
      \\[5mm]

    \item
      \(
        \lim_{x\to a} \left(f(x)g(x)\right)
        = \left(\lim_{x\to a} f(x)\right)\left(\lim_{x\to a} g(x)\right)
      \)
      \hfill Produto
      \\[5mm]

    \item
      \(
        \lim_{x\to a} \frac{f(x)}{g(x)}
        = \dfrac{\displaystyle \lim_{x\to a} f(x)}{\displaystyle \lim_{x\to a} g(x)}\quad
      \)
      desde que $\;\lim_{x\to a} f(x) \neq 0$
      \hfill Quociente
  \end{enumerate}
\end{frame}

%------------------------------------------------------------------------------%
\begin{frame}{Teorema do confronto}
  Se \(g(x) < j(x) < h(x)\)

  \pause
  para todo $x$ em um intervalo aberto contendo $a$,
  \pause
  exceto $x = a$,
  \pause
  e
  \[
    \lim_{x\to a} g(x) = \lim_{x\to a} h(x) = L
  \]
  \pause
  então
  \[
   \lim_{x\to a} f(x) = L
  \]
\end{frame}

%------------------------------------------------------------------------------%
\section{Lista Mínima}

%------------------------------------------------------------------------------%
\begin{frame}{Lista Mínima}
  Cálculo Vol. 1 do Thomas 12\textsuperscript{a} ed.
  \begin{enumerate}
    \item Revisar atentamente texto do Capítulo 2
  \end{enumerate}
  \vfill
  Atenção: A prova é baseada no livro, não nas apresentações
\end{frame}

%------------------------------------------------------------------------------%
\end{document}
%------------------------------------------------------------------------------%
